\chapter{Writing Formal Definitions and Proofs}
\label{ch:writing-foundations}

\section{The dependency order of a formal development}
\label{sec:writing-foundations-01}

A formal development should normally introduce material in the following
order:
\begin{enumerate}
  \item ambient sets and metalinguistic conventions;
  \item signature and arities;
  \item grammar of terms and formulas;
  \item free variable occurrences and substitution;
  \item structures and variable assignments;
  \item term evaluation and satisfaction;
  \item calculus and derivations;
  \item semantic and proof-theoretic properties; and
  \item connections between the calculus and semantics.
\end{enumerate}
An author may reorder material for a stated pedagogical reason, but no
definition may depend on an object that has not been defined.

\section{Definitions}
\label{sec:writing-foundations-02}

\begin{method}[Definition audit]
For every definition, answer the following questions in the definition or its
immediate context.
\begin{enumerate}
  \item What is the type or standing of the defined item?
  \item Over which previously defined objects does it range?
  \item Are all quantifiers explicit?
  \item Are all functions typed by domain and codomain?
  \item Are all arities and binding positions specified?
  \item Is the definition total, or is its domain of definition stated?
  \item If the definition is recursive, what is the well-founded measure or
        inductive grammar?
  \item If the definition uses a representative of an equivalence class, why is
        the result independent of that representative?
  \item If a new term is used, what previously unavailable distinction requires
        it?
\end{enumerate}
\end{method}

\begin{example}[Typed definition]
Let \(X\) be a set and let \(R\subseteq X\times X\).  Define the successor set
function \(R[-]:X\to\Pow(X)\) by
\[
  R[x]\defeq\set{y\in X:xRy}.
\]
The set \(X\), relation type, function type, argument, and value are all
explicit.
\end{example}

\begin{warning}[Circular definition]
A clause of the form ``a formula is an expression that is a formula according
to the usual rules'' does not define a set.  A reference to prior literature
may identify a convention, but a self-contained formal development must state
the grammar it uses.
\end{warning}

\section{Theorem statements}
\label{sec:writing-foundations-03}

A theorem statement should be readable independently of its proof.  It must
identify:
\begin{itemize}
  \item every object quantified over;
  \item every hypothesis, including finiteness, countability, consistency, and
        effectiveness assumptions;
  \item the exact language, calculus, and semantic class;
  \item the direction of each implication; and
  \item whether the result concerns formulas, sentences, theories, structures,
        or derivations.
\end{itemize}

The words ``clearly,'' ``obviously,'' and ``routine'' do not discharge a proof
obligation.  A short proof may cite a previous theorem and verify its
hypotheses; a missing proof may be marked as omitted, but it must not be
presented as established.

\begin{method}[Proof plan before proof]
Before writing details, record:
\begin{enumerate}
  \item the target statement with all quantifiers;
  \item the governing definition or theorem;
  \item the proof method;
  \item the induction parameter or minimal counterexample when applicable;
  \item the object to construct when the proof is existential; and
  \item the invariant maintained at each step.
\end{enumerate}
Delete the plan from a polished proof only if the proof remains easy to audit.
\end{method}

